\section{Markers and several scales of cuts}\label{sec:scales}

The main split will use a bounded family of possible cuts.  The clock
controls which tags can meet at a cut, but a suffix still has several
possible origins for its computation.  We therefore choose boundaries
whose interval products remain unchanged under suitable movements of the
origin.  Markers keep a boundary fixed outside periodic regions; inside
such a region we may move it by a multiple of a sufficiently large
period.  Several scales of movement ensure a uniform bound on the
number of ambiguous origins.

There are three distinct outputs. The marker rule recognizes regions where
letters determine a common periodic extension. The roster and scale
construction provides many cut positions of every required kind. Finally,
product transport and a descending count show that only a bounded number
of ordered pairs of those positions can give obstructed origins. We develop
the finite data in that order and record exactly which later choices may
still be enlarged.

All parameters in this section are finite and may depend on $M$ and
$\Sigma$; lengths, positions, and thresholds are integers.  We first
work on two-sided words, indexed by the integers.
On a finite argument, every window prescribed by a rule must be present.
An unavailable window causes rejection; in particular, it never removes
an origin from a universal or uniqueness test.  Section~\ref{sec:episodes}
will place all the required windows inside complete episodes.

\subsection{Local markers and periodic extensions}

A positive integer $a$ is a \emph{period} of a finite word if its letters
at distance $a$ agree whenever both positions belong to the word.

The marker-versus-local-periodicity viewpoint goes back to Krieger's marker
lemma~\cite[Lemma~2]{Krieger1982}; see also the finite clopen merging
formulation in Meyerovitch~\cite[Section~3, Lemmas~3.2--3.4]{Meyerovitch2025}.
We prove the finite-window rule needed here. The long-block hypothesis
$K>3d^2$ also prepares the overlap argument in the next lemma.

\begin{lemma}[Local markers]\label{lem:markers}
Let $d\geq 2$ and $K>3d^2$ be integers.  There is a
translation-equivariant rule marking integer positions of any two-sided
word such that distinct markers are at distance at least $d$.
The decision at $z$ uses only $[z-r,z+r)$, for a fixed finite integer
$r\geq K$.  If there is no marker in the inclusive integer interval
$[z-d,z+d]$, then the block $[z,z+K)$ has a positive period less than $d$.
\end{lemma}

\begin{proof}
List the finitely many length-$K$ words having no positive period below
$d$.  Process these words in a fixed order.  When processing a word,
mark all its starts that are not within distance less than $d$ of a
previously installed marker.  Two starts added at this same step cannot
be at distance $a<d$: their overlap would make $a$ a period of the
listed word.  Induction therefore proves separation of all markers.

Each step makes its decision from a length-$K$ block and the decisions
of earlier steps at offsets of magnitude less than $d$.  Induction over
the finite list gives a finite inspection radius; enlarge it to at
least $K$.  The construction commutes with translations.  Finally, a
listed word at $z$ was either marked or prevented by a marker within
distance less than $d$.  Thus marker absence on $[z-d,z+d]$ excludes
every listed word at $z$, as required.
\end{proof}

\begin{lemma}[Common periodic extensions]\label{lem:periodic-extension}
A block of length $K>3d^2$ having a positive period below $d$ determines
a unique two-sided periodic extension of period below $d$.  Two such
blocks whose starts differ by less than $K-d^2$ determine the same
extension.  For the marking rule of Lemma~\ref{lem:markers}, if there is
no marker in $[z-H,z+H]$ and $H\geq d$, the actual letters on
\[
       [z-H+d,\ z+H-d+K)
\]
belong to one such periodic extension.
\end{lemma}

\begin{proof}
A block of period $a$ extends by repeating its first $a$ letters.  Two
two-sided words of respective periods $a,b$ that agree on $ab$
consecutive positions agree everywhere: both have period $ab$.
Here $a,b<d$, so $ab<d^2$.  This proves both uniqueness and the assertion
about overlapping blocks.

For the last assertion, every start in the inclusive interval
$[z-H+d,z+H-d]$ has no marker within distance $d$.  Its length-$K$ block
therefore has a period below $d$.  Consecutive such blocks overlap
sufficiently, so their extensions agree.  Their union is the displayed
interval, on which the extension agrees with the actual letters.
\end{proof}

Whenever these lemmas are applied to a finite word, extend the word
arbitrarily in both directions.  If all prescribed inspection windows
are internal, the computed markers and periodic blocks are independent
of that extension.

\subsection{Algebraic data and the order of the buffers}

Before choosing cut positions, we fix the states and finite information
that each cut must transmit. Their cardinalities determine how many
candidate cut positions we need.

Put
\[
 U_0=V,\qquad U_{i+1}=\F_2^{\,U_i\times M}\quad (i=0,1).
\]
We denote the basis vector of $U_{i+1}$ indexed by $(v,c)$ by $[v,c]$.
Its work action is
\begin{equation}\label{eq:work-action}
          [v,c]\rho_i(d)=[v,cd]\qquad(d\in M).
\end{equation}
Thus the formal label $[v,c]$ retains an input state $v$ while its
second coordinate accumulates a prefix product. Once the later part of
an episode is known, a terminal linear map will send each label to a
basis label recording the completed update corresponding to $c$.
Lemma~\ref{lem:table-lift} constructs that map for every $c\in M$,
including values not realized by an actual prefix. The two larger spaces
will let us perform this encoding twice.

Choose $l,\beta$ from Lemma~\ref{lem:affine-line} for the action of $M$
on $V$.  For each $i=0,1$, choose $l_i,\beta_i$ from the same lemma for
the action $\rho_i$ on $U_{i+1}$.  These choices are made once and for all.
Thus $\beta:M^l\to V$ and $\beta_i:M^{l_i}\to U_{i+1}$.

An episode beginning at $s$ will have a later origin
\[
                         t=s+L.
\]
Before $t$ we will place two buffers in the order $1,0$.  Buffer $i$
has checkpoints
\[
 B_{i,0}<B_{i,1}<\cdots<B_{i,l_i}<a_i,
\]
where $a_i$ is after all its activities.  The entire buffer $1$,
including $a_1$, precedes buffer $0$.  The checkpoints are translates
of fixed offsets from $s$; we use $B_{i,j}$ and $a_i$ for their absolute
positions in the word.  The translations and $L$ will be fixed in
Section~\ref{sec:episodes}.

The \emph{kind} of a cut in buffer $i$ is a tuple
\begin{equation}\label{eq:buffer-kinds}
 (j,v,d_1,\ldots,d_{j-1}),\qquad
 1\leq j\leq l_i,\quad v\in U_{i+1},\quad d_f\in M.
\end{equation}
Let $\kappa_i$ be the number of these kinds.  A finite stencil
$\mathcal P_i$ will specify the cut offsets $p$ from $s$, with each
offset assigned a kind.  An offset of interval number $j$ must satisfy
$B_{i,j-1}<s+p<B_{i,j}$.  Write
\begin{equation}\label{eq:buffer-differences}
 h_i=\operatorname{diam}\mathcal P_i,\qquad
 \Delta_i=\mathcal P_i-\mathcal P_i.
\end{equation}
At an actual cut $s+p$, a suffix considering an offset $p'$ will infer
the origin $t+p-p'$.  We construct $\mathcal P_0$ here and
$\mathcal P_1$ in Section~\ref{sec:episodes}.

Once separated numerical offsets have been chosen, kinds can be
assigned in increasing interval-number order.  Boundaries between
successive interval numbers are then placed in the gaps between their
offsets; the extreme boundaries and $a_i$ can be added outside them.
Thus it suffices to produce separated offsets with sufficiently many
copies of each kind.

\subsection{The clock tags and a polynomial roster}

The first stencil must contain more copies of every kind than the total
number of wrong ordered origin pairs that can survive our tests. The
following polynomial roster leaves this growth margin: its number of slots
grows by a higher power of the field size than the nonclock obstruction
budget does. The exponents and other constants are fixed before that
field size is selected.

For the main segment set
\begin{equation}\label{eq:scale-counts}
             Z=l+1,\qquad k_0=2Z+1,\qquad R=2k_0-1.
\end{equation}
There will be $Z$ centers for its boundaries and $R$ scales for each
cut role.  Use the following tag set in Lemma~\ref{lem:clock}:
\begin{equation}\label{eq:clock-tags}
 \alpha=(j,r,v,d_1,\ldots,d_{j-1},X,Y),\qquad
 \begin{gathered}
 1\leq j\leq l,\quad 1\leq r\leq R,\\[-2pt]
 v,X,Y\in V,\quad d_f\in M.
 \end{gathered}
\end{equation}
Its ambiguity constant $C_{\rm cl}$ is determined by the tag set alone,
before a roster or its positions have been chosen.

Choose a prime $P>R$ large enough that
\begin{equation}\label{eq:stage0-constant}
 \left\lfloor\frac{P^{k_0}}{\kappa_0}\right\rfloor>C_0,
 \qquad
 C_0=C_{\rm cl}+2^R(3Z+1)^R P^{2Z}.
\end{equation}
Such primes exist because $k_0=2Z+1>2Z$.  The abstract roster consists
of all polynomials of degree less than $k_0$ over $\F_P$, including the
zero polynomial. This is the polynomial-evaluation construction underlying
Reed--Solomon codes~\cite{ReedSolomon1960}; we use only the elementary
root bound, which applies over the prime field chosen here.
Choose distinct $\xi_1,\ldots,\xi_R\in\F_P$ and give a
roster slot $p$ the digits
\[
                         x_r(p)=p(\xi_r).
\]
We identify field elements with $\{0,\ldots,P-1\}$ when using them as
digit indices.  A nonzero polynomial of degree less than $k_0$ has at
most $k_0-1$ roots, by successive division by linear factors.
Consequently any $k_0$ evaluations determine a slot, and two distinct
slots differ in at least
\begin{equation}\label{eq:roster-distance}
                         R-k_0+1=k_0
\end{equation}
digits.

Apply Lemma~\ref{lem:clock} to this roster, avoiding all primes at most
$N$.  Fix its modulus $n$, tag sets $I_\alpha,J_\alpha$, bad set $W$,
and roster residues $v_p\in\mathbb Z/n\mathbb Z$.  The prime divisors
of $n$ are all greater than $N$.  Notice that the modulus is chosen
after the roster size; enlarging it does not change $C_0$.

\subsection{Scale choices and the first stencil}

The abstract roster is now fixed, as are its required clock residues. We
next realize the slots at integer positions. Large separation at each
scale will dominate all smaller-digit changes, while a small residue
correction preserves the clock assignment.

Call a positive integer \emph{smooth} if all its prime divisors divide
$N$.  We choose thresholds $S_r<S'_r$, a positive step $\delta_r$,
a movement bound $R_r=n\delta_r$, and digit positions
$f_r(0)<\cdots<f_r(P-1)$, in increasing order of $r$.
For each $r$ perform the following choices:
\begin{enumerate}[label=(\roman*)]
\item Choose $S_r\geq1$, also requiring $S_r\geq S'_{r-1}$ if $r>1$.
Choose $\delta_r$ smooth and divisible by $N\sigma$ for every smooth
positive integer $\sigma\leq S_r$.  Then
\begin{equation}\label{eq:step-coprime}
                            \gcd(\delta_r,n)=1.
\end{equation}
\item If $h'_b$ denotes the diameter of the digit positions at scale
$b$, put
\[
                    J_r=n+\sum_{b<r}h'_b.
\]
Choose the digit positions so that all ordered differences of unequal
digits are distinct and separated from one another and from zero by
more than $10(J_r+R_r+1)$.  Denote their diameter by $h'_r$.
\item Choose
\begin{equation}\label{eq:upper-scale-threshold}
              S'_r>\max\{S_r,\,10(h'_r+J_r+R_r+1)\}.
\end{equation}
\end{enumerate}
The step in (i) may be the least common multiple of the finitely many
required $N\sigma$; it is smooth.  If $N=1$, the only smooth positive
integer is $1$, and $\delta_r=1$ is permissible.  The digit positions
in (ii) can be sufficiently scaled powers of two.  Indeed a positive
difference $2^a-2^b$, $a>b$, determines $b$ by its power-of-two
valuation and then determines $a$.  All nonzero ordered differences
are thus distinct, and scaling provides the required separation.

Assign slot $p$ the numerical offset
\begin{equation}\label{eq:slot-position}
                   \sum_{r=1}^R f_r(x_r(p))+\epsilon_p,
             \qquad 0\leq\epsilon_p<n.
\end{equation}
Choose $\epsilon_p$ to give residue $v_p$ modulo $n$.
For two slots whose highest differing digit is at scale $r$, the sum
of all smaller-coordinate differences and the correction difference
has absolute value less than $J_r$.  The difference at scale $r$
has absolute value greater than $10(J_r+R_r+1)$.  The resulting
positions are therefore distinct, and every two are separated by
more than $2$.

These positions, up to a common translation, form $\mathcal P_0$.
We henceforth identify a slot with its numerical offset when taking
differences.  Their difference residues are exactly the prescribed
$v_p-v_{p'}$, including after the common translation.  By
\eqref{eq:stage0-constant}, we can assign more than $C_0$ offsets to
each stage-$0$ kind.  Assign them in interval-number order and insert
the buffer boundaries as described above.  Translation of the complete
configuration later changes none of its difference properties.  In
particular,
\begin{equation}\label{eq:stage0-diameter}
                         h_0<n+\sum_r h'_r.
\end{equation}

\subsection{Inner controls and the main boundaries}

We have fixed the stencil and every movement range. We can now make the
inner markers sparse relative to all of them. In marker-free regions the
boundary rule must preserve a congruence class modulo a full monoid period;
choosing only a phase modulo the word period would not suffice.

Choose integers
\begin{equation}\label{eq:inner-marker-size}
 d_\circ>10\left(1+n+\sum_r h'_r+\sum_r R_r\right),
 \qquad K_\circ>3d_\circ^2,
\end{equation}
and fix the marking rule of Lemma~\ref{lem:markers} with these
parameters and radius $r_\circ\geq K_\circ$.  Choose an integer $H$ with
\begin{equation}\label{eq:inner-halfwidth}
             H>10\bigl((N+m+1)d_\circ+\max_r R_r+1\bigr).
\end{equation}
For an inner center $z$, let $u(z)$ be the first marker in
$[z-H,z+H]$, or $\bot$ if none exists.  If $u(z)\neq\bot$, put
$B(z)=u(z)$.

If $u(z)=\bot$, Lemma~\ref{lem:periodic-extension} supplies a periodic
extension of the block starting at $z$.  Let $d<d_\circ$ be its least
period, and let $\sigma$ be the largest smooth divisor of $d$.  Among
the length-$d$ period words starting at its different phases, choose
the lexicographically least, using a fixed order on $\Sigma$.  Its
starts form one class $x\bmod d$: equality of two rotations would
contradict minimality of $d$.  Define $y(z)$ to be the largest integer
at most $z$ in $x+\sigma\mathbb Z$, and let $B(z)$ be the least integer
at least $z$ satisfying
\begin{equation}\label{eq:boundary-congruences}
                  B(z)\equiv x\pmod d,\qquad
                  B(z)\equiv y(z)\pmod{\sigma N}.
\end{equation}
The definition is independent of the representative of $x$.  Moreover,
\[
 \gcd(d,\sigma N)=\sigma,\qquad
 \operatorname{lcm}(d,\sigma N)=dN,
\]
because $d/\sigma$ is coprime to $N$.  The two residues in
\eqref{eq:boundary-congruences} agree modulo $\sigma$, so they define
one class modulo $dN$.  In particular,
\begin{equation}\label{eq:periodic-boundary-range}
                            z\leq B(z)<z+dN.
\end{equation}
All these choices commute with translations of the letters.  They
do not use an external choice of position origin.  In both the present
and absent cases, $B(z)\in[z-H,z+H]$.

Choose an even integer
\begin{equation}\label{eq:center-spacing}
                    G>10(H+r_\circ+\max_r R_r+1).
\end{equation}
For the main segment set
\[
 o_j=(j+1)G,\quad z_j=t+o_j,\quad B_j=B(z_j)
       \quad(0\leq j\leq l),
 \qquad c_j=\frac{o_{j-1}+o_j}{2}\quad(1\leq j\leq l).
\]
At another proposed origin $\tau$, write
$z_j(\tau)=\tau+o_j$ and $B_j(\tau)=B(z_j(\tau))$; the notation
without an argument refers to the true origin $t$.
The main cut offsets from $t$ form the finite set
\begin{equation}\label{eq:main-cuts}
 \mathcal H=
 \{c_j+\delta_r a:1\leq j\leq l,\ 1\leq r\leq R,\ 0\leq a<n\}.
\end{equation}
A cut of role $(j,r)$ lies strictly between $B_{j-1}$ and $B_j$.
All inspection windows for boundaries through $j-1$ are before this
cut.  Conversely, for a suffix beginning at $k$ and any hypothesized
role $(j,r)$, a candidate origin $k-c_j-\delta_r a$ has its first
required center at
\begin{equation}\label{eq:suffix-first-center}
                       k+G/2-\delta_r a.
\end{equation}
Since $\delta_r a<R_r$, \eqref{eq:center-spacing} puts the inspections
for this center and all later centers after $k$.  This holds even when
the suffix role differs from the prefix role.  Complete placement of
these controls, including the later end controls, is deferred to
Lemma~\ref{lem:placement}.

\subsection{Usable scales and transport of products}

The boundary rule is defined at every origin. We next specify when a
whole range of shifted origins gives equal interval products. The
condition compares actual marker positions, including absence, and in
periodic regions controls the phase lattice used by the congruence rule.

\begin{definition}\label{def:usable}
A scale $r$ is \emph{usable at the origin $t$} if, at each of its $Z$
centers $z_j$, the following conditions hold:
\begin{enumerate}[label=(\roman*)]
\item The value $u(z_j+h)$ equals $u(z_j)$ for every integer $h$ with
$|h|\leq R_r$, including agreement on absence.
\item If $u(z_j)=\bot$, then either $\sigma\leq S_r$, or the inclusive
interval $[z_j-R_r,z_j+R_r]$ contains no point of
$x+\sigma\mathbb Z$.
\end{enumerate}
Here $d,\sigma,x$ refer to the periodic extension at the indicated
center.  All windows needed by these checks are mandatory.
\end{definition}

\begin{lemma}[Transport of boundary products]\label{lem:boundary-transport}
Suppose $r$ is usable at $t$, and $h$ is a multiple of $\delta_r$ with
$|h|\leq R_r$.  Write $B_j^{(h)}=B(z_j+h)$.  Then
\begin{equation}\label{eq:transport-intervals}
        T_{B_{j-1},B_j}=T_{B_{j-1}^{(h)},B_j^{(h)}}
                         \qquad(1\leq j\leq l).
\end{equation}
For a fixed endpoint $b\geq z_l+H$, the last transfer also satisfies
\begin{equation}\label{eq:transport-tail}
                         T_{B_l,b}=T_{B_l^{(h)},b}.
\end{equation}
These statements apply to finite words whenever the prescribed
windows and product intervals are available.  They hold for arbitrary
finite monoids, without a commutativity assumption.
\end{lemma}

\begin{proof}
The proof has three steps: the two boundaries stay in one class modulo
$dN$; the actual word contains enough periodic letters on both sides;
and monoid-power periodicity then compares the ordered products.

\emph{A common boundary class.}
Consider one center $z$. If its marker is present, usability leaves
the boundary unchanged.  If absent, the old and new periodic blocks
have the same extension by Lemma~\ref{lem:periodic-extension}:
$|h|\leq R_r<d_\circ$, and the blocks have length $K_\circ>3d_\circ^2$.
Thus the least period $d$, smooth divisor $\sigma$, and phase class
$x\bmod d$ are common.

If $\sigma\leq S_r$, the displacement $h$ is divisible by $\sigma N$.
Translation along the lattice gives $y(z+h)=y(z)+h$, so the second
residue in \eqref{eq:boundary-congruences} is unchanged.  Otherwise,
usability says that no point of that lattice lies between any of the
positions under consideration, so $y(z+h)=y(z)$.  In both cases the
two boundaries belong to one class modulo $dN$.  Writing them in
increasing order as $b_-\leq b_+$, we have
\begin{equation}\label{eq:boundary-displacement}
                            b_+-b_-=kdN\quad(k\geq0).
\end{equation}

\emph{Margins of actual periodic letters.}
By \eqref{eq:periodic-boundary-range}, both choices lie in
$[z-R_r,z+R_r+Nd)$.  Marker absence and
Lemma~\ref{lem:periodic-extension} give a common run of actual letters
containing $[z-H+d_\circ,z+H-d_\circ)$.  The inequality
\eqref{eq:inner-halfwidth} therefore implies that
\begin{equation}\label{eq:periodic-margins}
               [b_--md,\ b_++md)
                 \subset [z-H+d_\circ,z+H-d_\circ)
\end{equation}
lies in that run.  In particular both choices have the required
periodic margins on their two sides.

\emph{Comparing products one endpoint at a time.}
Let $c\in M$ be the product of one length-$d$ period beginning in
the phase of $b_-$.  For an incoming product from any fixed
$a\leq z-H$, comparison at $b_--md$ gives the two factorizations
\[
             T_{a,b_-}=A c^m,\qquad
             T_{a,b_+}=A c^{m+kN}.
\]
For an outgoing product to any fixed $b\geq z+H$, comparison at
$b_++md$ gives
\[
             T_{b_-,b}=c^{m+kN}C,\qquad
             T_{b_+,b}=c^m C.
\]
Here $A$ and $C$ are the products on the unchanged remaining portions.
Equation~\eqref{eq:monoid-periodicity} gives
$c^{m+kN}=c^m$.  The equalities require neither cancellation nor
reordering of factors.

The old and new boundary at each center lie within its unshifted
window $[z_j-H,z_j+H]$, by the same estimates.  Adjacent center
windows are disjoint and separated by more than their required
margins, since $G>10H$.  To compare an interval between successive
boundaries, move its left endpoint while fixing its right endpoint,
then move the right endpoint while fixing the new left endpoint.
Each fixed opposite endpoint is outside the relevant window, so the
local calculation applies.  This proves
\eqref{eq:transport-intervals}.  Its outgoing version at the last
center proves \eqref{eq:transport-tail}.
\end{proof}

\subsection{A uniform ambiguity bound for the first stencil}

Product transport handles an origin whenever at least one scale is usable.
We now bound all the pairs of stencil slots that can produce an origin
where this fails, also allowing a bad clock residue. Throughout this count
the anchor $q$ is fixed. Permitting it to vary with a candidate would not
give the translated-difference problem solved by the clock.

For a proposed origin $t$ and an anchor position $q$, define
\begin{equation}\label{eq:stage0-predicate}
 A_0(t,q):\qquad
     q-t\pmod n\in W
       \quad\text{or no scale is usable at }t.
\end{equation}
The anchor is held fixed when comparing different proposed origins.

\begin{proposition}[Inner ambiguity]\label{prop:inner-ambiguity}
On any two-sided word and for any integers $t,q$,
\begin{equation}\label{eq:inner-ambiguity}
 \#\{(p,p')\in\mathcal P_0^2:p\neq p',\
                         A_0(t+p-p',q)\}\leq C_0.
\end{equation}
The same bound holds on finite words whenever all the indicated
inspection windows are internal.
\end{proposition}

\begin{proof}
We count numerical difference tuples rather than choose digit pairs at
every stage. Nonzero differences recover the slots; zero differences
contribute only one value. This distinction prevents an unnecessary factor
of $P$ at each equal-digit coordinate.

\emph{Clock failures and recovery of the ordered pair.}
The clock part contributes at most $C_{\rm cl}$, because
\[
      q-(t+p-p')\equiv(q-t)+v_{p'}-v_p\pmod n.
\]
This is precisely the uniform translated-difference bound of
Lemma~\ref{lem:clock}.  It remains to count pairs for which no scale
is usable.

For an ordered pair of slots define its numerical digit differences
\[
                  D_r=f_r(x_r(p))-f_r(x_r(p')).
\]
The full tuple $(D_1,\ldots,D_R)$ determines a distinct ordered pair
$p\neq p'$.  Indeed, it has at least $k_0$ nonzero coordinates by
\eqref{eq:roster-distance}.  Each such coordinate determines both
ordered digits by the uniqueness of the digit differences.  These
$k_0$ evaluations determine each polynomial slot separately; unrevealed
equal digits at zero coordinates therefore introduce no additional pairs.

\emph{Exceptional scales.}
For a pair with no usable scale, look at its origin $t+p-p'$ and
declare a coordinate $r$ exceptional if
\begin{equation}\label{eq:exceptional-scale}
                S_r<\sigma\leq S'_r
\end{equation}
at at least one absent center.  At this fixed origin each center has
one value of $\sigma$, and the intervals $(S_r,S'_r]$ are disjoint.
There are therefore at most $Z$ exceptional coordinates.  We may sum
over all their possible sets, of which there are at most $2^R$.

\emph{Descending count with all smaller digits still free.}
Fix such a set of exceptions.  Count difference tuples in descending
coordinate order.  At an exceptional coordinate there are at most
$P^2$ possible differences.  At a nonexceptional coordinate $r$, fix
the differences at all larger coordinates.  For each center label
$j$, its possible positions are of the form
\begin{equation}\label{eq:center-digit-error}
 C_j+D_r+e,\qquad |e|\leq J_r,\qquad
 C_j=t+o_j+\sum_{b>r}D_b.
\end{equation}
The error includes the smaller-coordinate differences and
$\epsilon_p-\epsilon_{p'}$.  This description fixes only the larger
numerical differences; it does not assume that their underlying
digit pairs have already been determined.

Since scale $r$ fails to be usable, at some center label one of the
following occurs:
\begin{enumerate}[label=(\alph*)]
\item A marker lies within $R_r$ of one of the two search endpoints.
\item The search is absent at the center, $\sigma>S'_r$, and a point
of $x+\sigma\mathbb Z$ lies within $R_r$ of the center.
\end{enumerate}
Indeed a change in the first marker of a translated search interval
requires a marker to enter or leave through an endpoint.  If the
search does not change, failure of the second usability condition
requires absence, $\sigma>S_r$, and such a lattice point.
Nonexceptionality then strengthens $\sigma>S_r$ to $\sigma>S'_r$.

For (a), the range swept by a given search endpoint, enlarged by
$R_r$, has diameter at most
\begin{equation}\label{eq:swept-range}
                  2(h'_r+J_r+R_r)<d_\circ.
\end{equation}
It contains at most one marker.  A fixed marker can serve at most
one nonzero value of $D_r$: by \eqref{eq:center-digit-error}, serving
it restricts $D_r$ to an interval of length $2(J_r+R_r)$, whereas
distinct nonzero digit differences are more widely separated.
The two endpoints give at most two such values for each label.

For (b), all the centers under consideration for a fixed label lie
in an interval of diameter at most $2(h'_r+J_r)<d_\circ$.
Among those centers whose searches are absent, their length-$K_\circ$
blocks determine the same two-sided periodic extension, by
Lemma~\ref{lem:periodic-extension}: their start distance is less than
$d_\circ$, and $K_\circ-d_\circ>d_\circ^2$. Only absent candidates enter
this comparison. Present candidates between them are irrelevant, since
overlap directly compares any two of the absent candidates' seed blocks.
Thus, even as $D_r$, the lower digits, and the residue corrections vary, these absent candidates have one
common least period, smooth divisor, phase, and lattice.  This is
where the long local blocks are needed in the descending count.
If the common $\sigma$ does not exceed $S'_r$, none of these candidates
can realize (b).  Otherwise the enlarged center range has diameter
less than $\sigma$, by \eqref{eq:upper-scale-threshold}; it contains
at most one point of the common lattice.  That point accounts for
at most one nonzero difference $D_r$, by the same separation argument
as for a marker.

\emph{Summing the budgets.}
There are $Z$ center labels.  Allowing also the single difference
value zero, a nonexceptional coordinate therefore has at most
$3Z+1$ possibilities.  It is not necessary to specify the equal
digits when the difference is zero.  Summing the descending count
over exception sets bounds the number of full difference tuples by
\[
                  2^R P^{2Z}(3Z+1)^R.
\]
Since different ordered pairs give different tuples, this also bounds
the number of pairs.  Adding the clock contribution proves
\eqref{eq:inner-ambiguity} with the constant in
\eqref{eq:stage0-constant}.

For a finite word, take any two-sided extension.  Internal inspection
windows give the same decisions, so the bound just proved applies
unchanged.
\end{proof}

The order of all choices is worth recording.  The affine data and tag
set fix $C_{\rm cl}$; then $P$ fixes the roster; then the clock fixes
$n$; then the scale parameters and digit positions are chosen
successively.  Only after these choices do we choose the inner marker
parameters, $H$, and $G$.  Thus neither the ambiguity estimate nor the
number of copies per kind depends circularly on a later geometric
choice.  The remaining buffer, end controls, and translations can now
be added around these fixed finite arrangements.
